\documentclass[11pt]{article}

\usepackage[margin=1in]{geometry}
\usepackage[T1]{fontenc}
\usepackage{csquotes}
\usepackage{graphicx}
\usepackage{booktabs}

\usepackage[backend=biber,style=numeric,sorting=none,maxnames=3,minnames=1]{biblatex}
\addbibresource{bib/refs.bib}

\usepackage{hyperref}

\usepackage{bucket}

\hypersetup{
  colorlinks = true,
  linkcolor  = bucketblue,
  citecolor  = bucketblue,
  urlcolor   = bucketblue,
}

\newcommand{\emp}[1]{{\small[empirical: #1]}}
\newcommand{\bm}[1]{{\small[bm:#1]}}
\newcommand{\covrun}{data/coverage.json, 2026-11-25, python3 scripts/coverage.py}
\newcommand{\regrun}{data/regression.json, 2026-11-25, python3 scripts/regress.py}

\glossentry{coverage}{cognate coverage}%
  {https://github.com/gianyrox/bucket-foundation/blob/dev/lean/BucketMath/Learning/Coverage.lean}%
  {The token-weighted share of running text in a target language whose
   words a learner can link to a word in a language they hold, through a
   shared etymological ancestor or a close spelling; see \Cref{def:coverage}.}
\glossentry{held}{held set}%
  {https://github.com/gianyrox/bucket-foundation/blob/dev/lean/BucketMath/Learning/Coverage.lean}%
  {The set of languages a learner already reads, written $S$.}
\glossentry{closure}{ancestor closure}%
  {https://github.com/gianyrox/bucket-foundation/tree/dev/papers/cognate-coverage/scripts/coverage.py}%
  {The set of forms reachable from a word by following inheritance,
   derivation and reconstructed-root edges in the Wiktionary etymology
   graph, up to eight steps; see \Cref{def:closure}.}
\glossentry{fsi}{FSI class hours}%
  {https://www.state.gov/foreign-language-training}%
  {The class hours the US Foreign Service Institute reports for a native
   English speaker to reach professional working proficiency.}

\title{Cognate Coverage: How Much of a Target Language a Learner Already Reads}
\author{Bucket Foundation}
\date{2026-11-25}

\begin{document}
\maketitle

\begin{abstract}
We define \term{coverage}{cognate coverage} $C(S,t)$, the share of running
text in a target language $t$ whose words have a cognate or close form in
some language of a learner's \term{held}{held set} $S$, and prove in Lean
that it lies in $[0,1]$ and never drops when $S$ grows
\bm{BucketMath.Learning.Cognate.coverage\_mono}. We measure it over 23
targets and 27 held languages from a Wiktionary etymology and translation
index, over the analyzable words among each target's 5000 most frequent,
which carry 0.42 to 0.94 of their frequency mass \emp{\covrun}, weighted
by open word frequencies, counting only matches that
meet above Proto-Indo-European. The best single held language covers a
mean of 0.260 of target tokens, and the best set of four covers 0.351
\emp{\covrun}.
English coverage ranks \term{fsi}{FSI class hours} with Spearman $\rho=-0.669$, 95\% CI $[-0.865,-0.269]$, permutation $p=0.0003$, $n=23$; the linear fit explains 22\% of the variance and 72\% once the four 2200-hour languages are set aside [empirical: data/tests.json, 2026-11-25, python3 scripts/tests.py].
\end{abstract}

\section{Introduction}
\label{sec:intro}

A Spanish speaker who opens a Portuguese newspaper reads much of it on
the first day. A Spanish speaker who opens a Finnish one reads almost
none of it. Course designers and placement tests treat this gap as
folklore. We give it a definition, prove the two properties any
reasonable definition needs, and measure it.

The question has a set-valued input. Learners hold more than one
language, and a second held language can cover words the first misses: an
English speaker who also reads Latin recognizes Italian words English
lost. We compute $C(S,t)$ for every held set of size one to four over the
available languages and report the gain each added language brings.

\section{Related Work}
\label{sec:related}

Chiswick and Miller turned FSI-style learning scores into a quantitative
distance between English and other languages \cite{chiswick2005distance}.
Their measure is an outcome; ours is a lexical predictor of that outcome,
built from etymology and frequency. Token weights come from the
\texttt{wordfreq} package \cite{speer2018wordfreq}. The proofs use Lean 4
\cite{demoura2021lean4}.

\section{Preliminaries}
\label{sec:prelim}

Let $\mathcal{L}$ be a finite set of languages and $V_t$ the analyzable
vocabulary of a target $t$: the top 5000 words of $t$ by frequency that
have at least one entry in the translation table. Each $w\in V_t$ carries
a weight $f_t(w)$, its relative frequency in running text.

\begin{definition}[Ancestor closure]
\label{def:closure}
For a word $w$ of language $l$, the \term{closure}{ancestor closure}
$A(l,w)$ is the set of (language, form) pairs reachable from $(l,w)$ in at
most eight steps along Wiktionary etymology edges of kind inherited,
derived, form-of, or reconstructed root. For a target word the edges also
include borrowings; for a held-language word they do not.
\end{definition}

\begin{definition}[Recognition]
\label{def:recognition}
Let $E(w)$ be the English glosses of $w$ in the translation table, and
let $R_s(w)$ be the renderings of those glosses in a held language $s$,
restricted to the 10000 most frequent words of $s$ where \texttt{wordfreq}
covers $s$. Then $s$ \emph{recognizes} $w$, written $\kappa(s,w)=1$, when
some $u\in R_s(w)$ satisfies
\[
  \bigl(A(t,w)\cap A(s,u)\bigr)\setminus P\neq\emptyset
  \quad\text{or}\quad
  \frac{d(\phi(w),\phi(u))}{\max(|\phi(w)|,|\phi(u)|)}\le 0.25,
\]
where $\phi$ lowercases, romanizes from the table, and strips diacritics,
$P$ is the set of Proto-Indo-European forms, and $d$ is Levenshtein
distance. Exact matches need three letters and near matches need four.
Removing $P$ keeps matches a learner can see. The sensitivity variant
$\kappa_{\mathrm{PIE}}$ and $C_{\mathrm{PIE}}$ drop the $\setminus P$.
\end{definition}

\begin{definition}[Cognate coverage]
\label{def:coverage}
For a held set $S\subseteq\mathcal{L}$ with $t\notin S$,
\[
  C(S,t)=\frac{\sum_{w\in V_t} f_t(w)\,\max_{s\in S}\kappa(s,w)}
              {\sum_{w\in V_t} f_t(w)},
\]
with $\max_{s\in\emptyset}=0$.
\end{definition}

\begin{lemma}[Bounds]
\label{lem:bounds}
$0\le C(S,t)\le 1$ and $C(\emptyset,t)=0$
\bm{BucketMath.Learning.Cognate.coverage\_le\_one},
\bm{BucketMath.Learning.Cognate.coverage\_nonneg},
\bm{BucketMath.Learning.Cognate.coverage\_empty}.
\end{lemma}

\begin{theorem}[Monotonicity]
\label{thm:mono}
If $S\subseteq S'$ then $C(S,t)\le C(S',t)$
\bm{BucketMath.Learning.Cognate.coverage\_mono}.
\end{theorem}

\begin{proof}
Recognition is an existential over $S$, so a word recognized under $S$ is
recognized under $S'$ \bm{BucketMath.Learning.Cognate.recognized\_mono}.
The covered mass is a sum of nonnegative weights over the recognized
words, so it does not drop \bm{BucketMath.Learning.Cognate.coveredMass\_mono},
and dividing by the same nonnegative total keeps the order.
\end{proof}

The Lean model takes $\kappa$ as an arbitrary Boolean relation and the
text as a list of weighted words, so both results hold for any choice of
cognate test, including the one in \Cref{def:recognition}.

\section{Preregistration}
\label{sec:prereg}

Committed on 2026-11-25 before the rerun that produced every number in
\Cref{sec:results}. The FSI regression behind H1 had already been run once,
in commit 6575d95155, before this section was written, so H1 is
confirmatory only for the rerun with corrected FSI hours and
Proto-Indo-European matches excluded. H2 and H3 were specified before any
of their statistics were computed. The headline coverage $C$ counts a root match only
when the shared ancestor lies above Proto-Indo-European; $C_{\mathrm{PIE}}$
also counts matches that meet only at Proto-Indo-European and serves as a
sensitivity row. All resampling uses seed 20261125. Alpha is 0.05 for
every test.

\paragraph{H1, English coverage and FSI hours.}
H0: Spearman $\rho(C(\{\mathrm{en}\},t), h_t)\ge 0$. H1: $\rho<0$.
Test: one-sided permutation test, 10000 permutations of $h$ over the
targets FSI lists. Effect: $\rho$ with a 95\% percentile bootstrap
interval over languages, 10000 resamples, plus the OLS slope with its
bootstrap interval. Power: the smallest $|\rho|$ detectable at 80\% power,
by simulation from a bivariate normal at $n$. Secondary in the same family:
OLS of $\log h_t$ on $C$, slope tested by permutation. Holm correction
over the two.

\paragraph{H2, single-pair coverage.}
For each ordered pair $(s,t)$, H0: $C(\{s\},t)\le 0.05$. H1:
$C(\{s\},t)>0.05$. Test: one-sided $z$ test with the standard error from
1000 bootstrap resamples of $V_t$, each word carrying its weight.
Effect: $C(\{s\},t)-0.05$ with a 95\% bootstrap interval. Benjamini and
Hochberg at $q=0.05$ across all pairs. Power: the minimum detectable
excess $(z_{0.95}+z_{0.80})\cdot\mathrm{SE}$, reported as its median over
pairs.

\paragraph{H3, gain from a second held language.}
For each non-English target, split $V_t$ by frequency rank into odd and
even halves. On the odd half choose the $s\ne\mathrm{en}$ that maximizes
$C(\{\mathrm{en},s\},t)$. On the even half, H0: the gain
$C(\{\mathrm{en},s\},t)-C(\{\mathrm{en}\},t)\le 0.02$. H1: gain $>0.02$.
Test: one-sided bootstrap, 2000 resamples of the even half. Effect: the
gain with a 95\% interval. Holm correction across targets. Power: minimum
detectable gain as in H2.

\section{Method}
\label{sec:method}

The data are the Bucket photon index: a Wiktionary extraction with
5.9 million etymology edges and translation rows for every language in
the study \emp{\_intake/photons/roots.sqlite, 2026-11-25, sqlite3 roots.sqlite "select count(*) from etym"}.
Wiktionary content comes through Kaikki and is CC-BY-SA. The pgvector
index and the phonetic vectors were not used; the phonetic test is the
spelling distance of \Cref{def:recognition}.

Targets are the 23 study languages that \texttt{wordfreq} covers, plus
English. Held languages are those 23 plus Latin, Sanskrit, Thai and
English, 27 in all. For each target we compute $\kappa(s,\cdot)$ as a
Boolean vector over $V_t$ for each held $s$, then take unions for every
held set of size one to four that excludes $t$: $\binom{26}{4}=14950$ sets for a size of
four. We report the best set of each size, the mean over all sets of
each size, and two greedy chains: one free, and one that starts from
English.

For the regression we take FSI class hours for the 23 targets that FSI
lists: 600 for Category I, 750 for French and Spanish, 900 for German and Indonesian,
1100 for Category III, and 2200 for Arabic, Chinese, Japanese and Korean
\emp{data/fsi\_hours.json, 2026-11-25, Internet Archive snapshot 20241229203515 of state.gov/foreign-language-training}.
That page gives weeks per language; hours are weeks times 25 [empirical: data/fsi\_hours.json, 2026-11-25, Internet Archive snapshot 20241229203515]. It lists
Spanish at 30 weeks, so Spanish takes 750 hours. We fit
ordinary least squares of hours and of log hours on $C(\{\mathrm{en}\},t)$.

\section{Results}
\label{sec:results}

\Cref{fig:pairs} shows $C(\{s\},t)$ for every pair. The largest pairs sit
inside Romance: Portuguese covers 0.739 of Spanish tokens, Italian 0.645
of Spanish, Spanish 0.619 of Portuguese, and Latin 0.568 of Italian
\emp{\covrun}. Inside Germanic, Dutch covers 0.476 of German and English
covers 0.364 of Swedish. English covers under 0.05 of Czech, Greek,
Finnish, Japanese, Russian and Chinese tokens \emp{\covrun}.

\begin{figure}[t]
\centering
\includegraphics[width=0.85\textwidth]{figures/fig_pairs.png}
\caption{Single-language coverage $C(\{s\},t)$. Rows are targets, columns
held languages; the diagonal is undefined.}
\label{fig:pairs}
\end{figure}

\Cref{tab:sizes} and \Cref{fig:sizes} give coverage by set size, averaged
over the 23 non-English targets. The best set gains 0.057 from its second
language and 0.034 more from the next two. In the chain that starts from
English, the second language adds more than English did, 0.173 against
0.131, because it is a close relative of the target; the third and fourth
add 0.033 and 0.014 \emp{\covrun}.

\begin{table}[t]
\centering
\begin{tabular}{lcccc}
\toprule
$|S|$ & 1 & 2 & 3 & 4 \\
\midrule
mean over all sets & 0.043 & 0.079 & 0.109 & 0.135 \\
best set & 0.260 & 0.317 & 0.340 & 0.351 \\
greedy gain & 0.260 & 0.057 & 0.024 & 0.011 \\
gain, chain from English & 0.131 & 0.173 & 0.033 & 0.014 \\
\midrule
best single, $C_{\mathrm{PIE}}$ & 0.268 & & & \\
English alone, $C$ & 0.131 & & & \\
English alone, $C_{\mathrm{PIE}}$ & 0.151 & & & \\
\bottomrule
\end{tabular}
\caption{Coverage by held-set size, mean over 23 targets
\emp{\covrun}.}
\label{tab:sizes}
\end{table}

\begin{figure}[t]
\centering
\includegraphics[width=0.9\textwidth]{figures/fig_sizes.png}
\caption{$C(S,t)$ against $|S|$ for each target. Each curve rises or stays
flat, as \Cref{thm:mono} requires.}
\label{fig:sizes}
\end{figure}

\Cref{fig:fsi} plots FSI hours against English coverage. The linear fit
has slope $-2322$ hours per unit coverage, standard error 954, $R^2=0.220$,
$n=23$; on log hours $R^2=0.358$. Without the four 2200-hour languages,
$R^2=0.716$ on $n=19$ with slope $-1611$. With $C_{\mathrm{PIE}}$ in place
of $C$, $R^2=0.281$ \emp{\regrun}.

\subsection{Hypothesis Tests}
\label{sec:tests}

All results follow \Cref{sec:prereg} \emp{data/tests.json, 2026-11-25, python3 scripts/tests.py}.

\begin{table}[t]
\centering
\small
\begin{tabular}{llll}
\toprule
Hypothesis & Statistic and effect & $p$, corrected & Decision \\
\midrule
H1, $\rho<0$ & $\rho=-0.669$, CI $[-0.865,-0.269]$ & $0.0006$, Holm-adjusted & reject H0 \\
H1, log slope $<0$ & $-2.32$; OLS slope $-2322$ h, CI $[-3682,-965]$ & $0.0006$, Holm-adjusted & reject H0 \\
H2, $C(\{s\},t)>0.05$ & 48 of 624 pairs, 24 targets & BH-adjusted $\le 0.035$ & reject for 48 \\
H3, gain $>0.02$ & 12 of 23 targets & Holm-adjusted $\le 0.012$ & reject for 12 \\
\bottomrule
\end{tabular}
\caption{Preregistered tests. CI is the 95\% bootstrap interval.}
\label{tab:tests}
\end{table}

\textbf{H1.} The one-sided permutation test rejects H0 for both
statistics after Holm correction. At $n=23$ the smallest $|\rho|$ the test
detects with 80\% power is 0.55, and the observed $|\rho|$ exceeds it [empirical: data/tests.json, 2026-11-25, python3 scripts/tests.py].

\textbf{H2.} Benjamini and Hochberg reject H0 for 48 of 624 ordered pairs over 24 targets, English included.
The smallest rejected coverage is 0.076 and the largest unrejected is
0.220, a pair with wide bootstrap error. English as held language passes
for eight targets: German, Spanish, French, Indonesian, Italian, Dutch,
Portuguese and Swedish. The median minimum detectable excess over 0.05 is
0.0087.

\textbf{H3.} The second language chosen on the odd half adds more than
0.02 on the even half for 12 targets after Holm: Czech, German, Spanish,
Persian, French, Hindi, Italian, Dutch, Polish, Portuguese, Russian and
Swedish. Spanish gains most from Portuguese, 0.730, CI $[0.599,0.811]$.
H0 stands for Arabic, Greek, Finnish, Hebrew, Indonesian, Japanese,
Korean, Tamil, Turkish, Vietnamese and Chinese; Tamil comes closest, gain
0.031 from Sanskrit, uncorrected $p=0.063$.

\begin{figure}[t]
\centering
\includegraphics[width=0.55\textwidth]{figures/fig_fsi.png}
\caption{FSI class hours against $C(\{\mathrm{en}\},t)$, with the least
squares line.}
\label{fig:fsi}
\end{figure}

\section{Calibration}
\label{sec:calibration}

The bounds and monotonicity are derivations, proved in
\Cref{app:lean}. The measured values have no gold standard to calibrate
against; we checked them three ways. Every curve in \Cref{fig:sizes} is
nondecreasing, which a bug in the union step would break. Known close
pairs rank first (\Cref{fig:pairs}). Hand inspection of matched words
during development found two false-positive sources, which the method
now blocks: held-language loanwords that point back to the target through
English, such as Japanese \emph{za} for German \emph{die}, and rare
translation renderings. Blocking them lowered German coverage from
Japanese from 0.158 to under 0.03 on a development sample of the 300 most
frequent German words, run with \texttt{TOP\_N=300 TGT=de}; that sample
was not committed, and the current script with those settings reproduces
only the post-fix value.

\section{Limitations}
\label{sec:limits}

The regression has 23 points and FSI hours take five values, so the fit
measures rank agreement more than a curve. Hours also reflect grammar,
script and tones, which $C$ ignores; the 2200-hour languages all use
non-Latin scripts and sit far from the line.

$V_t$ holds 1325 to 3464 analyzable words per target, carrying 0.42 to
0.94 of the top-5000 token mass \emp{\covrun}. Many function words lack
translation rows and drop out.

The etymology graph is incomplete, so root matches undercount. Shared
roots that meet only at Proto-Indo-European are excluded from $C$; for
Russian that moves English coverage from 0.152 to 0.042 \emp{\covrun}.
Matches through other reconstructed stages, such as Proto-Germanic, stay
in. Chinese
characters shared by Chinese, Japanese and Korean are not modeled, so
Sino-Xenic coverage reads near zero. The spelling test credits chance
resemblance between unrelated words. Held-language borrowings are not
credited at all. The FSI hours come from a December 2024
archive snapshot; the live page refused automated fetches on 2026-11-25.

\printbibliography

\appendix

\section{Glossary}
\label{app:glossary}

\printglossaryentry{coverage}
\printglossaryentry{held}
\printglossaryentry{closure}
\printglossaryentry{fsi}

\section{Lean Listings}
\label{app:lean}

Source: \texttt{lean/BucketMath/Learning/Coverage.lean}, imported by
\texttt{BucketMath} and so by \texttt{BucketMathAll}. It defines
\texttt{Recognized}, \texttt{CoveredMass}, \texttt{TotalMass} and
\texttt{CognateCoverage} over $\mathbb{Q}$, with word weights as natural
numbers; the measured frequencies are rationals, and scaling them by a
common denominator turns them into such weights without changing $C$. It proves
\texttt{recognized\_mono}, \texttt{coveredMass\_mono},
\texttt{coveredMass\_le\_total}, \texttt{coverage\_mono},
\texttt{coverage\_le\_one}, \texttt{coverage\_nonneg} and
\texttt{coverage\_empty} with core Lean 4.33.1 only. Output of
\texttt{make lean}:

% voice-ignore-next 6
\begin{verbatim}
Built BucketMath.Learning.Coverage
Built BucketMath
Built BucketMathAll
Build completed successfully (31 jobs).
\end{verbatim}

\end{document}
